\chapter{Détection d'anomalies non supervisée}
\label{chap:anomalie}

\section{Motivation}
Les classifieurs supervisés (chapitre~\ref{chap:bench}) ne reconnaissent que les défauts
\emph{vus à l'entraînement}. Or une part importante des défaillances industrielles relève de
modes rares ou imprévus. Pour combler cet angle mort, PrognoSense intègre une couche de
détection d'anomalies \textbf{non supervisée}, entraînée uniquement sur des données saines :
tout écart au régime normal est signalé, qu'il corresponde ou non à un défaut catalogué.

\section{Auto-encodeur}
\subsection{Architecture}
L'auto-encodeur (PyTorch) est un réseau symétrique encodeur-décodeur qui apprend à
\emph{reconstruire} les vecteurs de features sains via un goulot d'étranglement de dimension
16 (espace latent). L'encodeur enchaîne des couches $d \to 256 \to 128 \to 64 \to 16$ avec
normalisation par lots (\emph{BatchNorm}), activations ReLU et \emph{Dropout} (0{,}2) ; le
décodeur est le miroir. La dimension d'entrée $d$ s'adapte au dataset (45 pour VBL, 15 pour
CWRU, 132 pour MF, 126 pour CMAPSS).

\subsection{Principe de détection}
Entraîné à minimiser l'erreur quadratique de reconstruction sur les seules données saines,
l'auto-encodeur reconstruit fidèlement un signal normal mais échoue sur un signal anormal.
L'\textbf{erreur de reconstruction}~\eqref{eq:recon} sert donc de score d'anomalie :
\begin{equation}
\varepsilon(x) = \tfrac{1}{d}\,\lVert x - \hat{x} \rVert_2^2
\label{eq:recon}
\end{equation}
Le \textbf{seuil de décision} est fixé au \textbf{95\textsuperscript{e} percentile} des
erreurs de reconstruction sur les données d'entraînement : 95\,\% des données saines sont
ainsi correctement reconnues comme normales. Les seuils réels obtenus sont reportés au
tableau~\ref{tab:ae}. Le score est ensuite normalisé sur $[0,100]$ par rapport aux
percentiles P5--P99 des erreurs d'entraînement.

\begin{table}[h]
\centering\small
\caption{Auto-encodeurs entraînés : dimension d'entrée et seuil (P95).}
\label{tab:ae}
\begin{tabular}{>{\bfseries}l c c}
\toprule
Dataset & Dimension d'entrée & Seuil P95 (erreur) \\
\midrule
VBL    & 45  & 0{,}1049 \\
CWRU   & 15  & 0{,}0493 \\
MF     & 132 & 0{,}5693 \\
CMAPSS & 126 & 0{,}3161 \\
\bottomrule
\end{tabular}
\end{table}

\begin{figure}[h]
\centering
\includegraphics[width=0.85\textwidth]{autoencoder_validation.png}
\caption{Validation de l'auto-encodeur : distribution des erreurs de reconstruction
sain vs défaut, et seuil de décision.}
\label{fig:ae}
\end{figure}

\section{Ensemble de détecteurs (vote par consensus)}
Pour fiabiliser la décision et \textbf{réduire les fausses alarmes}, un second dispositif
combine trois détecteurs d'anomalie aux biais complémentaires :
\begin{itemize}[leftmargin=1.5em]
  \item \textbf{Isolation Forest} : isole les points rares par partitionnement aléatoire,
        efficace en haute dimension ;
  \item \textbf{Local Outlier Factor} : compare la densité locale d'un point à celle de ses
        voisins ;
  \item \textbf{Elliptic Envelope} : suppose une distribution gaussienne et détecte les
        déviations de covariance.
\end{itemize}
Une anomalie n'est \textbf{confirmée que par vote majoritaire} : au moins deux des trois
détecteurs doivent s'accorder. Ce consensus diminue fortement le risque de fausse alerte
isolée. Pour produire des scores stables même sur un seul échantillon, chaque détecteur est
calibré contre une \textbf{référence de normalité} (percentiles P1--P99 de sa fonction de
décision sur les données saines), plutôt que normalisé au sein du lot courant.

\section{Apprentissage du baseline propre à chaque machine}
La détection non supervisée a été poussée plus loin avec le module \texttt{MachineBaseline}.
Plutôt que de comparer une machine à un dataset public générique, la plateforme apprend la
\textbf{signature saine spécifique de chaque machine} (statistiques multivariées : moyenne
et écart-type de chaque feature sur une fenêtre d'apprentissage). L'anomalie est alors
mesurée relativement à \emph{cette} machine :
\begin{equation}
z_i = \left|\frac{x_i - \mu_i}{\sigma_i}\right|, \qquad
s = \min\!\Big(100,\ 60\cdot \tfrac{1}{d}\textstyle\sum_i \mathbb{1}[z_i>3]
   + 40\cdot \min(1, \bar{z}/5)\Big)
\end{equation}
Ce score combine la proportion de features fortement déviantes ($z>3$) et l'ampleur moyenne
de la déviation. Validation : sur un signal sain le score reste $\approx 5\,\%$, sur un
signal dégradé il atteint $\approx 96\,\%$, avec identification des \emph{features
contributrices} (explicabilité). Cette approche --- apprendre le \og normal \fg{} de chaque
équipement sans aucun historique de panne --- est celle des solutions industrielles de
pointe, et constitue un apport majeur du projet (chapitre~\ref{chap:indus}).
